\documentclass[conference]{IEEEtran}
\IEEEoverridecommandlockouts

\usepackage{cite}
\usepackage{amsmath,amssymb,amsfonts}
\usepackage{algorithmic}
\usepackage{graphicx}
\usepackage{textcomp}
\usepackage{xcolor}
\usepackage{hyperref}
\usepackage{subfig}
\def\BibTeX{{\rm B\kern-.05em{\sc i\kern-.025em b}\kern-.08em
T\kern-.1667em\lower.7ex\hbox{E}\kern-.125emX}}

\begin{document}

    \title{Cross-Platform Quantum Compiler for Synthesizing Quantum Circuits Based on Arbitrary Hamiltonians \\
    \thanks{This work was supported in part by personal funding and community feedback.}
    }

    \author{\IEEEauthorblockN{Shoupu Wan}
    \IEEEauthorblockA{Independent Researcher \\
    Email: wanshoupu@gmail.com\thanks{Code available at: \url{https://github.com/wanshoupu/quantum-simulations}}}
    }

    \maketitle

    \begin{abstract}
        We present a quantum system software framework that bridges the gap between physics modeling and quantum software development.
        At the core of this framework is a versatile quantum compiler capable of decomposing arbitrary Hamiltonians into quantum circuits, represented using a platform-independent B-Tree intermediate language (IL).
        The B-Tree structure provides essential lineage information for circuit gate traceability and code verification.
        The IL offers a universal, hardware-agnostic representation that enables flexible translation and optimization across diverse quantum hardware backends.
        We demonstrate implementations that render the IL into executable quantum circuits targeting platforms such as Qiskit and Cirq.
        Additionally, we provide a transpiler that converts the IL into OpenQASM for broader compatibility.
    \end{abstract}

    \begin{IEEEkeywords}
        Quantum Computing, Quantum Compilation, Hamiltonian Decomposition, Cross-Platform, Circuit Synthesis, Solovay-Kitaev Theorem,
    \end{IEEEkeywords}


    \section{Relevance}\label{sec:relevance}
    Quantum simulation has demonstrated exponential advantages in applications such as chemistry, pharmaceuticals, and biological processes.
    The field of quantum software development, however, has a relatively high learning curve, requiring expertise and many years of training in physics, math, and software engineering,
    and this poses difficulty for many professionals to access quantum computing.
    Furthermore, the state-of-the-art choices of transpiler are either tied to a specific platform or lack the desired translation granularity.

    In this poster, we present a quantum circuit synthesizer that solves both problems.
    We introduce a design and implementation of quantum circuit synthesizer with good features to solve all the problems at hand.
    Our quantum circuit synthesizer features a cross-platform quantum compiler, built-in traceability and verifiability, and configurable gate control.
    First, the compiler accepts user-defined Hamiltonians as input and outputs B-tree data structure as IL\@.
    This gives the necessary lineage for each quantum gate through the parent-child relationship among the tree nodes.
    Second, the IL is platform-independent.
    It can be easily rendered into executables on any existing quantum computing platforms such as Qiskit and Cirq.
    It can also be transpiled into OpenQASM for broader compatibility.
    Third, the granularity of the IL, as an end result, is configurable.
    For example, for superconducting quantum hardware, Pauli gates $X$, $Y$, $Z$, may be precisely implemented via \emph{e.g.}, through virtual gates with zero-error and zero-latency
    whereas on trapped-ion platforms such as IonQ, Arbitrary axis-aligned rotations, e.g. $Rx\left(\phi\right)$, $Ry\left(\phi\right)$, and $Rz\left(\phi\right)$, are more fundamental.
    Our design puts the fine-tuning granularity control into the hands of users.
    Depending on the targeted hardware, the users can choose which types of gates are used in the end result.
    Another good feature is that the IL can be transpiled into OpenQASM\@.
    This expands the compatibility with other quantum computing platforms supporting QASM\@.
\end{document}
